% GENERATED FILE. Edit canonical lessons, then run npm run generate:books.
% canonical-source-hash: fnv1a64:4ffdc0228d054cc0
% canonical-lessons: KA-C194-sambala, KA-C194-grahaka, KA-C194-darji, KA-C194-manegelasa, KA-C194-meladhikari

\chapter{Salary, Customer, Tailor, Housework, Boss}
\label{ch:ka-salary-customer-tailor-housework-boss}
\hlchaptermodality{194}

\begin{chapteropening}
\textbf{By the end of this chapter:} \emph{I can say a salary, a customer, a tailor, housework and a boss.}

Complete the last lesson of chapter 194: I can say a salary, a customer, a tailor, housework and a boss.
\end{chapteropening}

\section[\texorpdfstring{\kn{ಸಂಬಳ} (sambaḷa)}{ಸಂಬಳ (sambaḷa)}]{\kn{ಸಂಬಳ} (sambaḷa) — a salary}

\label{lesson:KA-C194-sambala}



\begin{quote}
Before the new one: say the Kannada for a square, then the Kannada for a triangle.
\end{quote}

\subsection*{You'll want to know: \kn{ಸಂಬಳ}}
\textbf{\kn{ಸಂಬಳ}} — \emph{sambaḷa} — \textquotedblleft{}a salary\textquotedblright{}.

\begin{cousinweb}
Where \textbf{\kn{ಸಂಬಳ}} (\emph{sambaḷa}) comes from is not settled: DEDR does not record it. It may be Sanskrit \emph{śambala}, \textquotedblleft{}provisions for a journey\textquotedblright{} (CDIAL 12315), but CDIAL does not list the Kannada word.
\end{cousinweb}

\begin{grammarlens}[title={What you've built}]
One more word for this chapter.
\end{grammarlens}

\subsection*{Your turn}
\begin{itemize}
\raggedright
  \item \emph{Say it:} \emph{sambaḷa}
  \item \emph{Say it:} \emph{sambaḷa}, once more
  \item \emph{Say it:} say \emph{sambaḷa}
  \item \emph{Recall:} say the Kannada for a square, then the Kannada for a triangle, then say \emph{sambaḷa} again
\end{itemize}

\begin{tcolorbox}[breakable,colback=teal!4,colframe=teal!35!black,title={Before you move on}]
What does \kn{ಸಂಬಳ} mean? (\textquotedblleft{}A salary\textquotedblright{}.) Say it once more.
\end{tcolorbox}
\section[\texorpdfstring{\kn{ಗ್ರಾಹಕ} (grāhaka)}{ಗ್ರಾಹಕ (grāhaka)}]{\kn{ಗ್ರಾಹಕ} (grāhaka) — a customer}

\label{lesson:KA-C194-grahaka}



\begin{quote}
Before the new one: say the Kannada for a triangle, then the Kannada for a salary.
\end{quote}

\subsection*{You'll want to know: \kn{ಗ್ರಾಹಕ}}
\textbf{\kn{ಗ್ರಾಹಕ}} — \emph{grāhaka} — \textquotedblleft{}a customer\textquotedblright{}.

\begin{cousinweb}
\textbf{\kn{ಗ್ರಾಹಕ}} (\emph{grāhaka}) is Sanskrit \emph{grāhaka}, \textquotedblleft{}one who takes, a taker\textquotedblright{}, taken over whole: the customer is the one who takes the goods.
\end{cousinweb}

\begin{grammarlens}[title={What you've built}]
One more word for this chapter.
\end{grammarlens}

\subsection*{Your turn}
\begin{itemize}
\raggedright
  \item \emph{Say it:} \emph{grāhaka}
  \item \emph{Say it:} \emph{grāhaka}, once more
  \item \emph{Say it:} say \emph{grāhaka}
  \item \emph{Recall:} say the Kannada for a triangle, then the Kannada for a salary, then say \emph{grāhaka} again
\end{itemize}

\begin{tcolorbox}[breakable,colback=teal!4,colframe=teal!35!black,title={Before you move on}]
What does \kn{ಗ್ರಾಹಕ} mean? (\textquotedblleft{}A customer\textquotedblright{}.) Say it once more.
\end{tcolorbox}
\section[\texorpdfstring{\kn{ದರ್ಜಿ} (darji)}{ದರ್ಜಿ (darji)}]{\kn{ದರ್ಜಿ} (darji) — a tailor}

\label{lesson:KA-C194-darji}



\begin{quote}
Before the new one: say the Kannada for a salary, then the Kannada for a customer.
\end{quote}

\subsection*{You'll want to know: \kn{ದರ್ಜಿ}}
\textbf{\kn{ದರ್ಜಿ}} — \emph{darji} — \textquotedblleft{}a tailor\textquotedblright{}.

\begin{cousinweb}
\textbf{\kn{ದರ್ಜಿ}} (\emph{darji}) is a loan from Persian \emph{darzī}, \textquotedblleft{}tailor\textquotedblright{}, from \emph{darz}, \textquotedblleft{}seam\textquotedblright{}, which came through Hindustani.
\end{cousinweb}

\begin{grammarlens}[title={What you've built}]
One more word for this chapter.
\end{grammarlens}

\subsection*{Your turn}
\begin{itemize}
\raggedright
  \item \emph{Say it:} \emph{darji}
  \item \emph{Say it:} \emph{darji}, once more
  \item \emph{Say it:} say \emph{darji}
  \item \emph{Recall:} say the Kannada for a salary, then the Kannada for a customer, then say \emph{darji} again
\end{itemize}

\begin{tcolorbox}[breakable,colback=teal!4,colframe=teal!35!black,title={Before you move on}]
What does \kn{ದರ್ಜಿ} mean? (\textquotedblleft{}A tailor\textquotedblright{}.) Say it once more.
\end{tcolorbox}
\section[\texorpdfstring{\kn{ಮನೆಗೆಲಸ} (manegelasa)}{ಮನೆಗೆಲಸ (manegelasa)}]{\kn{ಮನೆಗೆಲಸ} (manegelasa) — housework}

\label{lesson:KA-C194-manegelasa}



\begin{quote}
Before the new one: say the Kannada for a customer, then the Kannada for a tailor.
\end{quote}

\subsection*{You'll want to know: \kn{ಮನೆಗೆಲಸ}}
\textbf{\kn{ಮನೆಗೆಲಸ}} — \emph{manegelasa} — \textquotedblleft{}housework\textquotedblright{}.

\begin{cousinweb}
\textbf{\kn{ಮನೆಗೆಲಸ}} (\emph{manegelasa}) is native Dravidian: \emph{mane}, \textquotedblleft{}house\textquotedblright{} (DEDR 4776), and \emph{kelasa}, \textquotedblleft{}work\textquotedblright{} (DEDR 1970), with the \emph{k-} softened to \emph{g-} inside the compound.
\end{cousinweb}

\begin{grammarlens}[title={What you've built}]
One more word for this chapter.
\end{grammarlens}

\subsection*{Your turn}
\begin{itemize}
\raggedright
  \item \emph{Say it:} \emph{manegelasa}
  \item \emph{Say it:} \emph{manegelasa}, once more
  \item \emph{Say it:} say \emph{manegelasa}
  \item \emph{Recall:} say the Kannada for a customer, then the Kannada for a tailor, then say \emph{manegelasa} again
\end{itemize}

\begin{tcolorbox}[breakable,colback=teal!4,colframe=teal!35!black,title={Before you move on}]
What does \kn{ಮನೆಗೆಲಸ} mean? (\textquotedblleft{}Housework\textquotedblright{}.) Say it once more.
\end{tcolorbox}
\section[\texorpdfstring{\kn{ಮೇಲಧಿಕಾರಿ} (mēladhikāri)}{ಮೇಲಧಿಕಾರಿ (mēladhikāri)}]{\kn{ಮೇಲಧಿಕಾರಿ} (mēladhikāri) — a boss, a superior}

\label{lesson:KA-C194-meladhikari}



\begin{quote}
Before the new one: say the Kannada for a tailor, then the Kannada for housework.
\end{quote}

\subsection*{You'll want to know: \kn{ಮೇಲಧಿಕಾರಿ}}
\textbf{\kn{ಮೇಲಧಿಕಾರಿ}} — \emph{mēladhikāri} — \textquotedblleft{}a boss, a superior\textquotedblright{}.

\begin{cousinweb}
\textbf{\kn{ಮೇಲಧಿಕಾರಿ}} (\emph{mēladhikāri}) is native \emph{mēl}, \textquotedblleft{}above\textquotedblright{} (DEDR 5086), and \emph{adhikāri}, \textquotedblleft{}an official\textquotedblright{}, from Sanskrit \emph{adhikārin}, \textquotedblleft{}one in charge\textquotedblright{}: \textquotedblleft{}the official above\textquotedblright{}.
\end{cousinweb}

\begin{grammarlens}[title={What you've built}]
One more word for this chapter.
\end{grammarlens}

\subsection*{Your turn}
\begin{itemize}
\raggedright
  \item \emph{Say it:} \emph{mēladhikāri}
  \item \emph{Say it:} \emph{mēladhikāri}, once more
  \item \emph{Say it:} say \emph{mēladhikāri}
  \item \emph{Recall:} say the Kannada for a tailor, then the Kannada for housework, then say \emph{mēladhikāri} again
\end{itemize}

\begin{tcolorbox}[breakable,colback=teal!4,colframe=teal!35!black,title={Before you move on}]
What does \kn{ಮೇಲಧಿಕಾರಿ} mean? (\textquotedblleft{}A boss, a superior\textquotedblright{}.) Say it once more.
\end{tcolorbox}
